\documentclass[11pt]{article}
\usepackage[a4paper,margin=25mm]{geometry}
\usepackage{fontspec}
\defaultfontfeatures{Ligatures=TeX}
\setmainfont{DejaVu Serif}
\setsansfont{DejaVu Sans}
\usepackage{polyglossia}
\setdefaultlanguage{ukrainian}
\setotherlanguage{english}
\newfontfamily\cyrillicfont{DejaVu Serif}[Script=Cyrillic]
\newfontfamily\cyrillicfontsf{DejaVu Sans}[Script=Cyrillic]
\usepackage{amsmath,amssymb,mathtools,bm}
\usepackage{booktabs,array,graphicx}
\usepackage{xcolor}
\usepackage{caption}
\usepackage{enumitem}
\usepackage[unicode,colorlinks,allcolors=blue!55!black]{hyperref}
\setlength{\emergencystretch}{2em}
\raggedbottom
\captionsetup{font=small,labelfont=bf}

\newcommand{\llbracket}{\mathopen{[\![}}
\newcommand{\rrbracket}{\mathclose{]\!]}}
\usepackage{framed,needspace}
\newcommand{\responseheading}{\textbf{Відповідь авторів.}\enspace}
\newcommand{\revisionheading}{\textbf{Внесені зміни.}\enspace}
\newenvironment{reviewcomment}{\def\FrameCommand{\textcolor{blue!45!black}{\vrule width 1.2pt}\hspace{9pt}}\MakeFramed{\advance\hsize-\width\FrameRestore}\noindent\textbf{Зауваження рецензента.}\par\smallskip\noindent}{\endMakeFramed}
\renewcommand{\thefootnote}{\arabic{footnote}}
\setlength{\parindent}{0pt}
\setlength{\parskip}{5pt plus 1pt}
\author{}
\date{}
\begin{document}

% R0001
\title{Відповіді рецензентам}
\maketitle
\vspace{-2em}
\begin{center}
Український переклад для наукового обговорення між співавторами.\par
Офіційні відповіді до журналу підготовлено англійською мовою.
\end{center}
\bigskip

% R0002
Дякуємо обом рецензентам за уважне прочитання рукопису та конструктивні зауваження. Рецензент~1 указав на кілька місць, де аналітичні міркування було викладено надто стисло. Незалежне відтворення підтвердило подані аналітичні вирази. Тому ми докладніше виклали проміжні етапи, уточнили прийнятий спосіб факторизації Вінера--Гопфа, розширили обговорення фізичної допустимості й додали кількісне порівняння двох конкурентних механізмів руйнування. Рецензент~2 рекомендував публікацію без конкретних вимог щодо змін.

% R0011
\textbf{Нумерація формул.} Номери у процитованих зауваженнях стосуються початкового варіанта, поданого до журналу. Нові проміжні формули змінили нумерацію в доопрацьованому рукописі. Початкові формули (1)–(5) зберегли свої номери. Початковим формулам (6)–(11) відповідають відповідно формули (9), (10), (11), (15), (16) і (17) доопрацьованого рукопису. У відповідях нижче за потреби вказано номери формул доопрацьованого рукопису.

% R0017
\pdfbookmark[1]{Рецензент 1}{reviewer-1}
\section*{Рецензент 1}

% R0019
\begin{reviewcomment}
У статті побудовано аналітичну модель для прогнозування умов, за яких наявна міжфазна тріщина зсуву, розташована в кутовій точці ламаної межі поділу двох пружних матеріалів, ініціює відгалужену тріщину з кутової точки всередину одного з матеріалів. Загалом виведення строгі й послідовні, і ця стаття заслуговує на публікацію в IJF. Однак моє головне зауваження полягає в тому, що в кількох місцях логічні переходи в міркуваннях дещо різкі, через що за виведенням важко стежити. Тому рекомендую авторам відповідно доопрацювати рукопис.
\end{reviewcomment}

% R0029
\responseheading Дякуємо рецензентові за позитивну оцінку й за те, що він указав, де додаткові проміжні етапи полегшили б розуміння виведення. Розділи 2–5 рукопису доопрацьовано у відповідь на наведені нижче докладні зауваження.

% R0034
\Needspace{7\baselineskip}
\pdfbookmark[2]{Зауваження 1. Формула (1)}{comment-1}
\subsection*{Зауваження 1. Формула (1)}

% R0036
\begin{reviewcomment}
У першому рядку вже зазначено, що \(\tau_{r\theta}=0\) для \(\theta\in\{0,\pi\}\), і незрозуміло, чому та сама умова знову з’являється в другому рядку. По-друге, я не бачу різниці між \(\llbracket\tau_{r\theta}\rrbracket=0\) і \(\tau_{r\theta}=0\) при \(\theta=\alpha\). Оскільки з першої, як видається, безпосередньо випливає друга, було б доречно пояснити фізичну або математичну різницю між ними. Крім того, оскільки автори обґрунтовують формулу (1) умовами симетрії, неперервності напружень і контакту без тертя, поетапне пояснення ролі кожної з цих умов у виведенні допомогло б читачам стежити за аналізом.
\end{reviewcomment}

% R0047
\responseheading Погоджуємося, що в початковому викладі три фізично різні умови було зведено до одного рядка. У доопрацьованій формулі (1) тепер окремо записано

% R0050
\begin{enumerate}
\item симетрію на \(\theta=0,\pi\), звідки \(\tau_{r\theta}=0\),
\item неперервність напружень на \(\theta=\alpha\), звідки \(\llbracket\sigma_\theta\rrbracket=0\) і \(\llbracket\tau_{r\theta}\rrbracket=0\),
\item контакт без тертя й розкриття берегів наявної міжфазної тріщини, звідки спільне дотичне напруження \(\tau_{r\theta}=0\) та умова \(\llbracket u_\theta\rrbracket=0\).
\end{enumerate}

% R0059
Також ми явно вказали, що \(\llbracket\tau_{r\theta}\rrbracket=0\) виражає рівність двох граничних дотичних напружень. Їх спільне значення з цієї умови не випливає. Воно дорівнює нулю за окремою умовою контакту без тертя.

% R0064
\revisionheading Розділ 2.2, формула (1) та абзац одразу після неї.

% R0067
\Needspace{7\baselineskip}
\pdfbookmark[2]{Зауваження 2. Формула (2)}{comment-2}
\subsection*{Зауваження 2. Формула (2)}

% R0069
\begin{reviewcomment}
Хоча я приблизно розумію фізичний зміст формули (2), у теперішньому вигляді її виклад читається доволі незграбно. Задля ясності його слід переписати.
\end{reviewcomment}

% R0073
\responseheading Погоджуємося. Формулу (2) переписано як умову, задану окремо на ділянках активної бісектриси. У доопрацьованому тексті тепер явно сказано, що на ділянці зони передруйнування \(0<r<d_i\) задано напруження, її продовження \(r>d_i\) залишається замкненим за симетрією, а бісектриса іншого матеріалу залишається лінією симетрії для всіх \(r>0\).

% R0079
\revisionheading Розділ 2.2, формула (2), а також вступний і пояснювальний текст до неї.

% R0081
\Needspace{7\baselineskip}
\pdfbookmark[2]{Зауваження 3. Формули (3) і (4)}{comment-3}
\subsection*{Зауваження 3. Формули (3) і (4)}

% R0083
\begin{reviewcomment}
Обґрунтуванню швидкості спадання \(o(r^{-1})\) у формулі (3), так само як і обґрунтуванню умови регулярності \(k_i=0\), варто присвятити по одному окремому реченню. По-перше, швидкість спадання в умові зрощування є саме тією умовою, за якої \(\Phi_i^+(p)\) залишається аналітичною в півплощині \(\operatorname{Re}p<0\), як того вимагає розділ 3. Вона не є внутрішньою властивістю зовнішнього поля. Явна вказівка на цей зв’язок зробила б наближення прозорим, а не неявним. Так само умова \(k_i=0\) у розділі 2.3 фактично є умовою розв’язності задачі з вільною межею. Явне зіставлення кількості невідомих та умов значно прояснило б, чому випливає формула (7). Нарешті, речення про те, що формула (4) є стандартною асимптотикою поля біля вершини в однорідному матеріалі, допомогло б читачам, які запитують, куди зник біматеріальний характер задачі.
\end{reviewcomment}

% R0097
\responseheading Погоджуємося й явно роз’яснили всі три пункти.

% R0099
По-перше, у доопрацьованому тексті вказано, що залишковий член \(o(r^{-1})\) виражає прийняту умову спадання коригувального поля після виділення провідного зовнішнього сингулярного члена. Її роль полягає в забезпеченні аналітичності зовнішньої трансформанти Мелліна при \(\operatorname{Re}p<0\). Також зазначаємо, що цієї умови достатньо, хоча вона й не мінімальна в математичному сенсі. У відкритій лівій півплощині вже було б достатньо слабшої оцінки \(O(r^{-1})\).

% R0106
По-друге, тепер пояснюємо, чому формула (4) має стандартний для моди~I в однорідному матеріалі вигляд із квадратним коренем. Вершина зони передруйнування \(r=d_i>0\) лежить усередині матеріалу \(i\) поза кутовою точкою межі поділу, тому провідний член її асимптотичного поля відповідає однорідному матеріалу.

% R0111
По-третє, у доопрацьованому рукописі \(d_i\) явно визначено як невідому кінцеву точку однопараметричної задачі з вільною межею. За заданої кінцевої точки розв’язок трансформованої задачі загалом дає ненульовий локальний коефіцієнт інтенсивності напружень \(k_i\). Умова регулярності \(k_i=0\) є скалярною умовою розв’язності й визначає рівноважне значення \(d_i\).

% R0117
Ми також уточнили знакову домовленість для переміщення на березі, використаного в перетворенні. Для трансформанти градієнта переміщення додатним вважаємо напрямок розкриття.

% R0121
\revisionheading Розділи 2.1–2.3 і визначення \(\Phi_i^-\) в розділі 3.1.

% R0124
\Needspace{7\baselineskip}
\pdfbookmark[2]{Зауваження 4. Розділи 3.1–3.2}{comment-4}
\subsection*{Зауваження 4. Розділи 3.1–3.2}

% R0126
\begin{reviewcomment}
Два проміжні етапи побудови за методом Вінера--Гопфа варто показати явно й чітко вказати, що саме факторизується. По-перше, фраза «Виключення трансформант внутрішнього поля дає функціональне рівняння» згортає все виведення рівняння (5). Стислий виклад цього етапу або його перенесення до додатка дав би читачеві змогу простежити міркування без самостійного відтворення, а також показав би, як отримано визначники додатка A. По-друге, безпосередній розгляд визначників показує, що ядро \(G_i(p)\) має полюси при \(p=\pm1\) і \(p=-1-\lambda\) та усувну особливу точку при \(p=0\). Проте у формулах (6)–(11) множники обчислюють саме в цих точках. Імовірно, тривіальні сингулярні множники вилучають із ядра й явно враховують у розв’язку, а факторизують регуляризоване ядро. Явна вказівка на це усунула б єдину неоднозначність, що залишилася в побудові. Так само тотожність для гамма-функціональних множників і виклад міркування за теоремою Ліувілля в трьох рядках зробили б розділ 3.2 самодостатнім.
\end{reviewcomment}

% R0143
\responseheading Погоджуємося, що ці проміжні етапи слід показати.

% R0145
У доопрацьованому розділі 3.1 тепер наведено стисле виведення рівняння (5). Для компоненти Мелліна функцію напружень Ері записано у вигляді \(r^{1-p}f(\theta)\) з кутовими модами частот \(1-p\) і \(1+p\). Застосування умови нульового дотичного напруження на бісектрисах секторів разом з умовами на переміщення й напруження на межі поділу та виключення амплітуд на межі поділу дають

% R0152
\[
T_i(p)=-\frac{D_i(p)}{D(p)}\Phi_i^-(p).
\]

% R0156
Поєднання цього співвідношення з \(D_i/D=\tan(\pi p)G_i\) і трансформантою заданого напруження на активній ділянці дає рівняння (5).

% R0160
У розділі 3.2 тепер явно вказано, що за допомогою інтеграла Коші факторизуємо додатний граничний символ \(G_i(it)\) з нульовим індексом на уявній осі. Для спрямованого вгору контуру співвідношення Сохоцького--Племеля дає факторизацію у вигляді відношення множників

% R0164
\[
G_i(it)=\frac{G_i^+(it)}{G_i^-(it)}.
\]

% R0168
Для парного ядра маємо \(G_i^-(p)=1/G_i^+(-p)\) у відповідних точках двох півплощин. Отже, полюси мероморфного продовження поза контуром факторизації не створюють особливих точок відповідного плюс-фактора Коші при \(-1\) або \(-1-\lambda\).

% R0173
Також ми додали тотожність для гамма-функціональних множників

% R0175
\[
K^+(p)K^-(p)=p\cot(\pi p)
\]

% R0179
і показали розділення на ліву й праву аналітичні частини з виділенням полюсних доданків перед застосуванням теореми Ліувілля. Ці доповнення явно показують перехід від рівняння (5) до доопрацьованого рівняння (9), тобто початкового рівняння (6).

% R0183
Незалежне відтворення розв’язку задачі для двох секторів у просторі трансформант і факторизації підтвердило правильність формул (5) та (6) у поданому вигляді. Під час доопрацювання додано пропущені етапи виведення зі збереженням самих формул, які тепер мають номери (5) і (9).

% R0188
\revisionheading Розділи 3.1 та 3.2, зокрема формули (5)–(9) доопрацьованого рукопису.

% R0190
\Needspace{7\baselineskip}
\pdfbookmark[2]{Зауваження 5. Розділ 3.4}{comment-5}
\subsection*{Зауваження 5. Розділ 3.4}

% R0192
\begin{reviewcomment}
Якби було показано обчислення \(\Phi_i^-(p)\) при \(p=0\) та \(p=1\), читач міг би безпосередньо перевірити формули (9)–(11). У переході від (6)–(7) до (9)–(10) є кілька неявних етапів, які варто виписати. Відтворюючи ці етапи, я отримую вираз для \(\delta_i\), в якому \(G_i(0)\) стоїть у чисельнику, тоді як у формулі (9) він стоїть у знаменнику. Я пов’язую цю розбіжність із домовленістю щодо множників в усувних особливих точках, яку обговорено в моїх попередніх зауваженнях. Явний проміжний етап дав би змогу однозначно встановити, де стоїть цей множник.
\end{reviewcomment}

% R0201
\responseheading Дякуємо рецензентові за те, що він указав на цю неоднозначність. Ми незалежно вивели два моменти трансформанти. Початкова формула (9), тепер формула (15), правильна в поданому вигляді, але прийнятий спосіб факторизації було викладено недостатньо явно, щоб місце множника \(G_i(0)\) було очевидним.

% R0206
Після накладання умови \(k_i=0\) визначимо

% R0208
\[
A_i=\frac{\sigma_iK^+(-1)}{G_i^+(-1)}.
\]

% R0212
Тоді доопрацьоване рівняння (9), початкове рівняння (6), набуває вигляду

% R0214
\[
\Phi_i^-(p)=
\frac{\lambda A_iK^-(p)G_i^-(p)}
{(p+1)(p+1+\lambda)}.
\]

% R0220
За факторизації у вигляді відношення множників, установленої в доопрацьованому розділі 3.2, маємо

% R0222
\[
G_i^-(0)=\frac{1}{\sqrt{G_i(0)}},
\qquad
K^-(0)=\frac{1}{\sqrt{\pi}},
\]

% R0228
і тому

% R0230
\[
\Phi_i^-(0)=
\frac{\sigma_i\lambda K^+(-1)}
{(1+\lambda)G_i^+(-1)\sqrt{\pi G_i(0)}}.
\]

% R0236
Отже, \(\sqrt{G_i(0)}\) стоїть у знаменнику доопрацьованої формули (15), тобто початкової формули (9).

% R0238
Аналогічно,

% R0240
\[
K^-(1)=K^+(-1)=\frac{2}{\sqrt{\pi}},
\qquad
G_i^-(1)=\frac{1}{G_i^+(-1)},
\]

% R0246
звідки

% R0248
\[
\Phi_i^-(1)=
\frac{2\lambda\sigma_i}
{\pi(\lambda+2)[G_i^+(-1)]^2}.
\]

% R0254
Це безпосередньо відтворює коефіцієнти у формулах для роботи й енергетичного параметра в доопрацьованій формулі (16), тобто початковій формулі (10).

% R0257
Під час незалежної числової перевірки для \(E_1/E_2=0.5\), \(\nu_1=\nu_2=0.3\), \(\alpha=45^\circ\) та \(\sigma'=-0.5\) отримали

% R0261
\[
d_1/l=0.147434361719,\qquad
\delta'_1=0.145787032394,\qquad
J'_1=0.083620387223,
\]

% R0267
що узгоджується з наявними результатами регресійної перевірки в MATLAB.

% R0269
\revisionheading Розділ 3.4, доопрацьовані формули (12)–(17). Усі коефіцієнти початкових формул (9)–(11) збережено.

% R0272
\Needspace{7\baselineskip}
\pdfbookmark[2]{Зауваження 6. Розділ 4}{comment-6}
\subsection*{Зауваження 6. Розділ 4}

% R0274
\begin{reviewcomment}
Числовий аналіз підтверджує заявлені тенденції, але два фізичні механізми вибору й твердження про допустимість лише однієї зони залишаються неперевіреними. По-перше, твердження, що умови допустимості \(Cg_2<0\) та \(CQ_i>0\) «можуть виконуватися одночасно лише для одного з двох матеріалів», явно обмежено «наборами параметрів, розглянутими нижче». Оскільки весь прогноз напрямку відгалуження в статті спирається на цю взаємовиключність, її слід або довести для всієї допустимої області параметрів, або явно подати як спостереження, отримане в результаті числового аналізу, з чітким зазначенням його області чинності. По-друге, у статті визначено конкурентний механізм ініціювання, але не порівняно критичних рівнів навантаження двох механізмів. Потрібне кількісне порівняння.
\end{reviewcomment}

% R0286
\responseheading Погоджуємося з обома зауваженнями й відповідно доопрацювали рукопис.

% R0289
Щодо першого зауваження, ми не узагальнювали результат про допустимість зони передруйнування лише в одному з матеріалів до аналітичної теореми. Ми явно зазначили, що це числовий результат, і окреслили область його чинності. Також провели значно ширшу перевірку, ніж у поданому варіанті.

% R0294
З виразу для \(q_i\) в додатку A випливає, що множник поза функцією чисельника \(S_i\) є спільним для обох розглянутих матеріалів. Отже,

% R0297
\[
\operatorname{sgn}(Q_1Q_2)
=
\operatorname{sgn}(S_1S_2).
\]

% R0303
Для одночасної допустимості за фіксованого знака навантаження потрібно, щоб виконувалося \(Q_1Q_2>0\), і тому \(S_1S_2>0\). Далі проведено детермінований перебір 92~800 невироджених випадків у діапазоні

% R0307
\[
10^{-3}\le E_1/E_2\le10^3,
\]

% R0311
Для відношення модулів узято 25 значень за логарифмічною шкалою, для кожного матеріалу — вісім незалежно варійованих значень коефіцієнта Пуассона від \(-0.95\) до \(0.49\), а також \(\alpha=3^\circ,6^\circ,\ldots,177^\circ\), за винятком виродженої границі прямолінійної межі поділу. Випадків із \(S_1S_2>0\) не знайдено.

% R0316
У другому дослідженні випадкової вибірки з 20~000 додаткових комбінацій із фіксованим початковим станом генератора відношення модулів вибирали рівномірно в логарифмічному масштабі в тому самому діапазоні, а коефіцієнти Пуассона й кут — неперервно в діапазонах

% R0320
\[
-0.99<\nu_i<0.499,
\qquad
0.2^\circ<\alpha<179.8^\circ.
\]

% R0326
Знову не виявлено жодного випадку, який міг би задовольняти умову одночасної допустимості. У двох випадках вибірки, найближчих до \(S_1S_2=0\), повторні обчислення з 50 десятковими цифрами точності підтвердили від’ємний знак. Співвідношення для взаємної заміни матеріалів також перевірено незалежно на 1~000 додаткових випадково вибраних комбінацій.

% R0331
Відповідно, у доопрацьованому рукописі зазначено лише, що допустимість зони передруйнування лише в одному з матеріалів підтверджено чисельно в дослідженій області. Загальної теореми єдиності не стверджуємо. Обидві умови допустимості й далі застосовуємо незалежно до кожного з двох матеріалів.

% R0336
Щодо другого зауваження, ми додали кількісне порівняння критичних навантажень. У позначеннях цієї статті розв’язок Назаренка й Кіпніса для коефіцієнта інтенсивності напружень за модою~II у вершині гілки має вигляд

% R0340
\[
K_{\mathrm{II}}
=
g_1k\,C\,l^{\lambda_0+1/2},
\]

% R0346
де

% R0348
\[
k=
\frac{
2\sqrt{2}(1+\kappa_2e)\Gamma(\lambda_0+3/2)
}{
[1+\kappa_1+(1+\kappa_2)e]\,
\Gamma(\lambda_0+2)G^+(-1-\lambda_0)
}.
\]

% R0358
Відповідний поріг для вершини гілки дорівнює

% R0360
\[
|C_{\mathrm{tip}}|
=
\frac{K_{\mathrm{IIc}}}
{|g_1k|\,l^{\lambda_0+1/2}}.
\]

% R0367
Поріг відгалуження з кутової точки, отриманий за \(\delta_i=\delta_{i\mathrm{c}}\) або \(J_i=J_{i\mathrm{c}}\), дорівнює

% R0371
\[
|C_{\mathrm{corner}}|
=
\sigma_i l^{-\lambda_0}|\sigma'_c|.
\]

% R0377
Визначення нормованої міжфазної тріщиностійкості за модою~II

% R0379
\[
\mathcal R_{\mathrm{II}}
=
\frac{K_{\mathrm{IIc}}}{\sigma_i\sqrt l},
\]

% R0385
дає

% R0387
\[
\boxed{
\frac{|C_{\mathrm{corner}}|}{|C_{\mathrm{tip}}|}
=
\frac{|g_1k|\,|\sigma'_c|}
{\mathcal R_{\mathrm{II}}}
}.
\]

% R0396
Отже, порога відгалуження з кутової точки досягаємо першим за умови

% R0398
\[
|g_1k|\,|\sigma'_c|<\mathcal R_{\mathrm{II}},
\]

% R0402
а за протилежної нерівності першим досягаємо порога поширення від вершин гілок уздовж межі поділу.

% R0405
Для базового варіанта \(E_1/E_2=0.5\), \(\nu_1=\nu_2=0.3\) та \(\alpha=45^\circ\) отримуємо

% R0409
\[
g_1k=-0.6794157816.
\]

% R0413
За нормованого критичного розкриття

% R0415
\[
\Delta_1=
\frac{E_1\delta_{1\mathrm{c}}}
{4(1-\nu_1^2)\sigma_1l},
\]

% R0421
межа між областями реалізації конкурентних механізмів набуває вигляду

% R0423
\[
\mathcal R_{\mathrm{II}}^*
=
1.1064\,\Delta_1^{0.6132}.
\]

% R0429
Рівнозначно, визначення

% R0431
\[
\mathcal J_1=
\frac{E_1J_{1\mathrm{c}}}
{4(1-\nu_1^2)\sigma_1^2l},
\]

% R0437
дає

% R0439
\[
\mathcal R_{\mathrm{II}}^*
=
1.5558\,\mathcal J_1^{0.6132}.
\]

% R0445
Обидва подання визначають ту саму межу за \(J_{1\mathrm{c}}=\chi_1\sigma_1\delta_{1\mathrm{c}}\).

% R0448
Для цього порівняння не потрібно задавати довільних значень тріщиностійкості. Для однозначного визначення порядку реалізації механізмів у конкретній системі матеріалів потрібні незалежно задані значення \(K_{\mathrm{IIc}}\) і однієї з величин \(\delta_{i\mathrm{c}}\) або \(J_{i\mathrm{c}}\).

% R0453
\revisionheading У вступі та розділі 2.1 тепер явно пов’язано припущення про задану стаціонарну тріщину з порівнянням критичних навантажень після розв’язання задачі. У розділі 2.2 твердження про взаємовиключність умов допустимості визначено як спостереження, отримане в результаті числового аналізу. У розділі 4 наведено розширений аналіз допустимості та його область параметрів. У розділі 5.2 додано формулу для \(K_{\mathrm{II}}\) у вершині гілки (доопрацьована формула (18)), два критичні навантаження, безрозмірне відношення критичних навантажень (доопрацьована формула (19)) та межу між областями реалізації конкурентних механізмів для конфігурації з кутом \(45^\circ\). Відповідні висновки також доопрацьовано.

% R0463
\pdfbookmark[1]{Рецензент 2}{reviewer-2}
\section*{Рецензент 2}

% R0465
\begin{reviewcomment}
Роботу присвячено дослідженню локальної зони передруйнування поблизу кутової точки V-подібної міжфазної тріщини зсуву та використанню її характеристик для прогнозування ініціювання відгалуження в один із матеріалів. У маломасштабному наближенні припускають, що довжина зони в матеріалі значно менша за довжину гілки тріщини. Визначено, у якому матеріалі може існувати фізично допустима зона. Перетворення Мелліна, зрощування асимптотик і метод Вінера--Гопфа використано для одержання замкнених аналітичних співвідношень для умов фізичної допустимості, рівноважної довжини й розкриття зони, роботи сил опору відриву, зонального енергетичного параметра та показника залишкової сингулярності в кутовій точці. Розкриття та енергетичний параметр використано відповідно в деформаційному й енергетичному критеріях ініціювання відгалуженої тріщини.

% R0465-P2
Рукопис добре написано, він містить нові цікаві результати, які можуть бути корисними для наукової спільноти, і його можна рекомендувати до публікації.
\end{reviewcomment}

% R0481
\responseheading Дякуємо рецензентові~2 за позитивну оцінку й уважне прочитання рукопису. Конкретних виправлень рецензент не вимагав.

\end{document}
